% Empirical Strategy --- Why Banks Open Branches in the Digital Era
% Last edited: 2026-05-12

\section{Empirical Strategy}
\label{sec:identification}

The empirical question asks how each of the three hypothesized channels --- local deposit contestability, the bank's cost-side edge against incumbents, and the bank's existing non-branch lending presence --- shapes the bank's decision to open a branch in a given ZIP within its candidate set. We need a specification that reads these channel coefficients off variation \emph{within} the bank--year (so that bank-level strategic shifts and aggregate bank shocks are absorbed) and \emph{within} the state--year (so that state-level regulatory and macroeconomic shocks are absorbed), after conditioning on the local demand variables described in Section~\ref{sec:data:sources}. The residual identifying variation is then cross-sectional, across the ZIPs in the bank's candidate set in a given year, holding the bank's overall expansion intensity and the state's overall environment fixed.

For each bank--ZIP--year cell, we estimate a linear probability model of the form
\begin{equation}
\mathbf{1}\{\mathrm{open}_{b,z,t}\}
\;=\;
\alpha_{b,t} \;+\; \gamma_{s(z),t} \;+\;
\beta\,\beta^{\mathrm{dep}}_{b,z,t}
\;+\;
\lambda\,C_{b,z,t-1}
\;+\;
\delta'\,X_{z,t-1}
\;+\;
\theta'\,Z^{\mathrm{footprint}}_{b,z,t-1}
\;+\;
\varepsilon_{b,z,t},
\label{eq:main}
\end{equation}
where the dependent variable is the opening indicator defined in Section~\ref{sec:data:sources}, $\alpha_{b,t}$ is a bank--year fixed effect, $\gamma_{s(z),t}$ is a state--year fixed effect (for the state $s(z)$ that contains ZIP $z$), $\beta^{\mathrm{dep}}_{b,z,t}$ is the imputed bank--ZIP--year deposit beta, $C_{b,z,t-1}$ is one of the channel variables (rate-paying advantage, interest-expense undercut share, mortgage presence, or CRA small-business presence), $X_{z,t-1}$ is the ZIP-level control vector measured at $t-1$, and $Z^{\mathrm{footprint}}_{b,z,t-1}$ collects the three existing-footprint controls (the bank's deposit share, log of one plus its own branch count, and the same-state indicator). We estimate the model separately for each of the three asset-size buckets defined in Section~\ref{sec:data:sources}, because the channel coefficients are not constrained to be equal across the size distribution and the patterns in Section~\ref{sec:data:desc} suggest they are not. To read the contestability--edge interaction the channel literature predicts, we additionally estimate the specification with $\beta^{\mathrm{dep}}_{b,z,t}$ entering interactively with the channel variable.

The identifying assumption is that, conditional on the bank--year and state--year fixed effects, the ZIP-level control vector, and the existing-footprint controls, the residual variation in $\beta^{\mathrm{dep}}_{b,z,t}$ and in the channel variables is orthogonal to other unobserved drivers of the bank's marginal ZIP choice. The bank--year fixed effect $\alpha_{b,t}$ absorbs every component of the bank's choice that is constant across the ZIPs the bank evaluates in year $t$ --- its overall expansion intensity, balance-sheet position, strategic posture, and aggregate exposure to the rate cycle --- so the contestability and channel coefficients are read off how the bank chooses between ZIPs inside its candidate set rather than off how aggressive a branch builder it is in the aggregate. The state--year fixed effect $\gamma_{s(z),t}$ absorbs every component of the local environment that is common to ZIPs in the same state in the same year, including state-level macroeconomic and regulatory shifts. The ZIP-level control vector strips out the demand-side variation that an entry model must net out before reading bank--market channels: local credit and deposit growth, demographics, density, and the indicators for absence of an incumbent or for in-state location. The three existing-footprint controls separate the marginal-entry decision from the lateral-expansion decision: a bank that already has a deposit relationship in a ZIP, that already operates a branch in the ZIP at $t-1$, or that already operates in the ZIP's state is a different decision-maker from a bank entering the ZIP without any of those.

Four potential threats to this reading are worth naming. First, reverse causality at the bank--market level --- a bank that has already decided to enter a ZIP at $t$ might have begun building up the lending presence that we use as a cross-sell channel at $t-1$. We mitigate this by measuring every channel and every ZIP-level control at $t-1$, and by treating the cross-sell channel as a binary indicator for any non-branch lending rather than the bank's actual lending volume, so that a single mortgage origination satisfies the indicator and the timing-driven reverse channel cannot scale up the coefficient. Second, the candidate-set restriction itself conditions on a prior decision: the bank has revealed willingness to operate in the CBSA. Our coefficients identify what determines the choice of ZIPs \emph{within} that CBSA footprint, not the CBSA-level choice. We treat this as a feature, not a bug --- the paper's question is about the marginal-market decision conditional on revealed metropolitan presence --- and the size-stratified estimation in particular keeps the within-CBSA comparison concentrated among banks for which the marginal-market choice is non-trivial. Third, the bank--market deposit beta is itself an imputed measure constructed from a within-cycle first stage. If the first-stage projection is mis-specified, $\beta^{\mathrm{dep}}_{b,z,t}$ is measured with error and $\beta$ in equation~(\ref{eq:main}) is biased toward zero; we let the first-stage coefficients shift across the three rate cycles (2004--2006, 2016--2019, 2022--2023) to accommodate time-varying pass-through, and we report the first-stage by cycle in Section~\ref{sec:results}. Fourth, bank-by-year fixed effects do not absorb a bank's \emph{market-specific} strategic plan that varies across its candidate ZIPs in a given year. Such a plan is observationally indistinguishable from the channel variables to the extent it loads on the same observables; we read this as part of the constructive interpretation of the channel coefficients --- "what makes a bank choose ZIP $z$ over ZIP $z'$ within the same CBSA" --- rather than as a separate confounder, but it does mean the coefficients should be read as conditional associations net of the design, not as outputs of a randomized experiment.

We cluster standard errors two-way at the bank and ZIP levels. Clustering at the bank level allows arbitrary correlation across all ZIPs and years for a given bank, which is the natural threat under bank-specific strategic shocks. Clustering at the ZIP level allows arbitrary correlation across all banks evaluating the same ZIP, which is the natural threat under ZIP-specific shocks correlated across banks (e.g., a local zoning change that makes branch entry mechanically easier or harder). The two-way clustering structure is conservative in the directions equation~(\ref{eq:main}) is most exposed to.
